\chapter{Imparare dagli esempi: il dataset}\label{cap:p1-dati}

% Capitolo della Parte I --- il dataset spiegato in modo semplice.

Per insegnare allo studente (la rete neurale) a riconoscere il rischio, serve
materiale su cui esercitarsi: degli \textbf{esempi}. Nel progetto questi
esempi si chiamano \textbf{dataset}.

\section{Cosa contiene il dataset}\label{sec:p1-dati-cosa}

Il dataset è una semplice raccolta di ``schede paziente''. Ogni scheda ha due
parti:

\begin{itemize}
  \item \textbf{I numeri} (i sei parametri vitali misurati);
  \item \textbf{La risposta} (il rischio giusto: basso, medio o alto), scritta
        dall'esperto di cui abbiamo parlato nel capitolo precedente.
\end{itemize}

È esattamente come un mazzo di \textbf{flashcard}: da una parte i dati del
paziente, dall'altra la risposta corretta. Lo studente le usa per imparare.

\section{Perché i dati sono ``finti''}\label{sec:p1-dati-sintetici}

I pazienti del dataset non sono persone vere: sono generati artificialmente
dal computer, ma in modo da essere il più possibile realistici. Tre buoni
motivi:

\begin{itemize}
  \item \textbf{Privacy}: non servono dati sanitari veri, quindi non ci sono
        problemi di riservatezza.
  \item \textbf{Controllo}: possiamo decidere quanti esempi generare per ogni
        tipo di situazione, così lo studente vede un po' di tutto.
  \item \textbf{Ripetibilità}: con gli stessi numeri iniziali (il ``seed'') il
        computer rigenera sempre esattamente gli stessi dati. Stessi dati,
        stessi risultati.
\end{itemize}

Nel progetto il dataset contiene 70\,000 schede così costruite.

\section{Tre mazzi diversi: prova, simulazione ed esame}\label{sec:p1-dati-split}

Non si usano tutte e 70\,000 le schede nello stesso modo. Si dividono in tre
mazzi, ciascuno con un ruolo:

\begin{itemize}
  \item \textbf{Allenamento (train)}: le schede su cui lo studente studia e
        impara. Sono la maggior parte (40\,000).
  \item \textbf{Verifica (validation)}: un mazzo usato per controllare, durante
        lo studio, se lo studente sta migliorando o sta solo imparando a
        memoria. Servono 10\,000 schede.
  \item \textbf{Esame (test)}: un mazzo tenuto da parte e guardato una sola
        volta alla fine, per misurare quanto lo studente sa davvero su casi
        mai visti. Sono 20\,000 schede.
\end{itemize}

L'importante è che l'esame non venga mai ``sbirciato'' prima: altrimenti le
misurazioni finali non sarebbero oneste.

\begin{esempio}{Una flashcard del progetto}
Una scheda del dataset potrebbe essere:
\begin{quote}
numeri: pressione 120/80, battiti 75, temperatura 36,8, ossigeno 98\%,
glicemia 95 \\
risposta: \textbf{basso}
\end{quote}
Lo studente vede i numeri e deve imparare ad arrivare da solo alla risposta
``basso''.
\end{esempio}
